\chapter{A direct all-tests join for the five-root A-unbounded subchart}
\label{component:p0d4eedd6c7a4}
\noindent\textbf{Component }\path{p0d4eedd6c7a4}\par
{\small\noindent Source: \path{R13/advances/weak_A/five_root_remote.tex}\par}
{\small\noindent SHA256: \path{0d4eedd6c7a4226bd797b288d93901cec2bcdae6b3ffe11bc5e98a60609914b4}\par}
{\small\noindent Inventory use: WR1-WR7\par}
\medskip
\section{A direct all-tests join for the five-root A-unbounded subchart}
\label{p0d4eedd6c7a4:sec:r9-weak-A-remote}

\paragraph{Revision 10 scope update.}
Section~\ref{p296e824cb97a:sec:r10-weak-boundary} extends this admission to $z/B$
positive infinitesimal with all other source hypotheses retained.
The unit-ratio restriction below records the original revision-9 theorem.

This section supplies a new direct join for exactly the source
hypotheses (WA1). It does not import the strong-II remote theorem
or any other weak-C2 remote admission. All source physicality,
strict-gap and contact hypotheses from the local theorem remain
in force. No membership in the residual complement is needed;
applications to the original residual problem retain that intersection
without asserting that it is nonempty.

\begin{theorem}
Under the source hypotheses of the five-root local theorem, including
the positive-unit assumptions on $z/B$ and $1-r$, the entire original
parent is tensor SOS for every independent projective testing angle
pair and every real or complex projective state. Literal zero and
infinite half-angles and all signed charts are included.
This is an all-tests source-subdomain admission; it does not assert
nonemptiness of its intersection with the previous residual complement.
\end{theorem}

\paragraph{The limiting source is explicit and all four finite weights are strict.}
Put
\[
 \rho=\operatorname{st}(z/B)>0,\qquad
 \tau=\operatorname{st}(y/B)>0,\qquad
 r_0=\operatorname{st}(r)\in(0,1).
\]
The original unit gaps give
$g_0=1-\rho(1-r_0)>0$ and
$j_0=1-1/\tau>0$, hence $\tau>1$.
Exact reduction of the unchanged source formulas gives
\[
 \begin{gathered}
 n_0=\frac{1-r_0}{2r_0(1+\rho)},\qquad
 m_0=\frac{\rho^2(1-r_0)}{2g_0(1+\rho)},\\
 C_0=\frac{\tau(1+\rho r_0)-\rho(1-r_0)}{2(1+\rho)},\qquad
 D_0=\frac{1+\rho r_0}{2(1+\rho)(\tau-1)}.             
 \end{gathered}
\tag{WR1}
\]
Every displayed coefficient is strictly positive. For $C_0$ the
numerator equals $\tau g_0+\rho(\tau-1)+\rho r_0>0$.
In particular $n$ is a positive bounded unit on this subchart;
the stronger fact is used at one remote endpoint. The formulas
are joint regular reductions under $B\to0$, $H/B\to0$, with
the other source ratios tending to their residues; they do not
compare the two infinitesimal rates. Also
$\alpha\in\operatorname{Inf}$ and $A^{-1/2}\in\operatorname{Inf}$.

Normalize the hard state by $\widehat Y=\sqrt A\,Y$ in each
remote chart. This is a pure-source state change, and no parent
term is removed. On bounded testing charts its matrix has bounded
pure-source polynomial coefficients; all couplings of
$\widehat Y$ to $W,X,Z$ have zero residue. Its hard diagonal
residue is the positive angle factor multiplying $AY^2$.
It remains to check the complete three-state residue on $Y=0$.

\paragraph{Finite angles outside the local corner.}
For finite real angle residues $u_0,v_0$, write $U=u_0^2$,
$V=v_0^2$, $D_u=1+U$, $D_v=1+V$. Let $\mathsf P_3$ be the
matrix of the following complete residue of $F/4$ on $Y=0$:
\[
 \begin{split}
 F_3={}&n_0D_uD_vX^2+m_0(v_0W-X-u_0Z)^2
                    +C_0D_vU W^2\\
 &+(D_0+\tfrac12)(u_0W+u_0v_0X-v_0Z)^2
       +\tfrac12D_v(X+u_0Z)^2-\tfrac12D_uVZ^2.
 \end{split}                                           \tag{WR2}
\]
No hard-state relaxation term remains in this residue, because
the normalized $\widehat Y$ cross row is infinitesimal. Define
the strictly positive real constants
\[
 a_*=\frac{D_0m_0(2n_0+1)}2,\quad
 c_* =\frac{C_0(2m_0+1)(2D_0+2n_0+1)}4,\quad
 d_* =\frac{C_0D_0(2n_0+1)}2.
\]
The exact residue determinant is
\[
 \det\mathsf P_3=D_uD_v
 \left\{a_*\left(V-\frac\tau\rho U\right)^2
                         +c_*U^2V+d_*UV^2\right\}.    \tag{WR3}
\]
The checker obtains this directly from the full three-state
polynomial and (WR1), not from a numerical table. It is strictly
positive whenever $U+V>0$, including either axis. The restriction
of (WR2) to $Z=0$ is positive definite on $(W,X)$: its positive
$n_0,m_0,C_0$ event squares already have no common nonzero kernel
when $U+V>0$, and the core restriction is nonnegative. The real
Schur criterion and (WR3) therefore give
$\mathsf P_3\succ0$. Together with the hard $\widehat Y$ pivot,
the entire four-state residue is strictly positive definite.

For independent bounded testing angles with these residues,
expand the normalized full matrix around the real pair
$(u_0,v_0)$. Its nonconstant monomials in $u-u_0,v-v_0$ are
pure-testing infinitesimals. The finite separated Gram construction
(WA10) applies verbatim. If $u_0=v_0=0$, use the local five-root
theorem instead. Thus every bounded testing pair is covered.

\paragraph{One infinite angle.}
For $u$ unbounded use $p=1/u$ and the full homogeneous parent
$p^2F(1/p,v)$; at literal $u=\infty$ use its polynomial value
at $p=0$. If $v$ is bounded, the complete three-state residue on
$Y=0$ is
\[
 C_0(1+v_0^2)W^2+n_0(1+v_0^2)X^2
       +(D_0+\tfrac12)(W+v_0X)^2+(m_0+\tfrac12)Z^2.
                                                               \tag{WR4}
\]
It is positive definite for every finite real $v_0$, including
$v_0=0$. The hard $\widehat Y$ pivot is $1+v_0^2>0$.
For $v$ unbounded and $u$ bounded, put $q=1/v$ and use the
polynomial parent $q^2F(u,1/q)$. Its complete three-state residue is
\[
 (m_0+C_0u_0^2)W^2+(n_0+\tfrac12)(1+u_0^2)X^2
                                      +D_0(u_0X-Z)^2, \tag{WR5}
\]
again positive definite, with hard pivot $1+u_0^2$.
Both identities are verified directly against the original parent.
The same bounded separated Gram lift applies in the infinitesimal
variables $p,v-v_0$ or $u-u_0,q$. This includes literal infinity
and literal zero of the other half-angle.

\paragraph{Both infinite angles: keep the exact null state.}
Put $p=1/u$, $q=1/v$, and work with the full polynomial parent
$\widehat F=p^2q^2F(1/p,1/q)$. Introduce the five roots
\[
 (W,X,\widehat Y,P,Q),\qquad
 \widehat Y=\sqrt A\,Y,\quad P=pZ,\quad Q=qZ.
\]
Write $Y=A^{-1/2}\widehat Y$ in the following exact quadratic
identity, interpreted as a polynomial in these five roots:
\[
 \begin{split}
 \widehat F/4={}&n(1+p^2)(1+q^2)X^2
   +m(pW-pqX+Y-Q)^2+C(1+q^2)(W-pY)^2\\
 &+(1+p^2)(1+q^2)\widehat Y^2
   +D(qW+X-pqY-P)^2+\alpha(1+p^2)(Y-Q)^2\\
 &+\tfrac12\left\{(qW+X-pqY-P)^2+(1+p^2)(Y-Q)^2
       +p^2(1+q^2)X^2\right.\\
 &\hspace{39mm}\left.+2(1+q^2)XP-(1+q^2)P^2\right\}.
 \end{split}                                           \tag{WR6}
\]
The checker verifies this as the complete original four-state
parent after $P=pZ$, $Q=qZ$, not only to leading order.
All matrix coefficients in (WR6) are bounded source polynomials
in $p,q$. Its residue, in the printed root order, is
\[
 \begin{pmatrix}
 C_0&0&0&0&0\\
 0&n_0+D_0+1/2&0&-D_0&0\\
 0&0&1&0&0\\
 0&-D_0&0&D_0&0\\
 0&0&0&0&m_0+1/2
 \end{pmatrix}\succ0,                                 \tag{WR7}
\]
because the nontrivial determinant is $D_0(n_0+1/2)>0$.
The separated lift (WA10) therefore gives a finite tensor SOS
for arbitrary independently infinitesimal $p,q$, of either sign
or zero. At $p=q=0$ both $P,Q$ vanish, respecting the original
exact null direction $W=X=Y=0$; no positive $Z$ pivot is invented.

\paragraph{Coverage and restoration.}
Every finite element of an ordered extension is either bounded or
unbounded in absolute value. A bounded testing pair is covered by
(WA1)--(WA10) if both residues vanish, and by (WR2)--(WR3)
otherwise. A single unbounded coordinate is covered by (WR4) or
(WR5); two unbounded coordinates are covered by (WR6)--(WR7).
Literal projective endpoints use the polynomial chart values with
$p=0$ or $q=0$. For nonzero reciprocal coordinates, restoring
$F$ multiplies the established identity by the pure-testing
square $p^{-2}$, $q^{-2}$, or $(pq)^{-2}$; at literal infinity
no division is performed. All four homogeneous state coordinates
were retained in every chart, and splitting real and imaginary
state parts gives the complex-state statement. Thus the join
exhausts all independent testing types without a source/testing
rate assumption or an imported strong-face theorem.

\begin{corollary}[The polynomial obstruction family is nevertheless admitted]
The entire two-parameter physical family (SO1),
\[
 B=\varepsilon,\quad z=\varepsilon/2,\quad y=2\varepsilon,
 \quad r=9/10,\quad H=\varepsilon\eta,
 \quad L=\varepsilon(1-\eta),
\]
with independent positive infinitesimals $\varepsilon,\eta$, is
all-tests tensor SOS for the original full parent.
\end{corollary}
\begin{proof}
It has $z/B=1/2$, $1-r=1/10$, $H/y=\eta/2$ infinitesimal,
and $y/B,L/B,r,G/B,j$ positive bounded units. Its global
physicality and all original source signs were proved in the
obstruction module. It therefore satisfies every hypothesis of
the present all-tests theorem.
\end{proof}
The strict $K_1<0$ and (JG13) failures and the nonexistence of a
global polynomial SOS of its scalar graph remain correct.
The local tensor roots use pure-testing radicals as in (WA10),
so these facts are compatible. The whole family is now a proved
admission and is removed from the residual complement. No claim
that any other residual intersection is nonempty is implied.
